\begin{thebibliography}{36}
\bibitem[1]{Ber90} V. G. Berkovich – Spectral theory and analytic geometry over non-Archimedean fields, Math. Surveys and Monographs, vol. 33, American Mathematical Society, Providence, RI, 1990.
\bibitem[2]{ber} V. G. Berkovich – “Vanishing cycles for formal schemes”, Invent. Math. 115 (1994), no. 3, p. 539–571.
\bibitem[3]{bo1} S. Bosch – “Eine bemerkenswerte Eigenschaft der formellen Fasern affinoider Räume”, Math. Ann. 229 (1977), no. 1, p. 25–45.
\bibitem[4]{che} V. I. Chernousov – “The Hasse principle for groups of type E8”, Dokl. Akad. Nauk SSSR 306 (1989), no. 5, p. 1059–1063.
\bibitem[5]{merkchern} V. I. Chernousov \& A. Merkurjev – “R-equivalence and special unitary groups”, J. Algebra 209 (1998), no. 1, p. 175–198.
\bibitem[6]{plati} V. I. Chernousov \& V. P. Platonov – “The rationality problem for semisimple group varieties”, J. reine angew. Math. 504 (1998), p. 1–28.
\bibitem[7]{cthhkps} J.-L. Colliot-Thélène, D. Harbater, J. Hartmann, D. Krashen, R. Parimala \& V. Suresh – “Localglobal principles for constant reductive groups over semi-global fields”, Michigan Math. J. 72 (2022), p. 77–144.
\bibitem[8]{ctps} J.-L. Colliot-Thélène, R. Parimala \& V. Suresh – “Patching and local-global principles for homogeneous spaces over function fields of p-adic curves”, Comment. Math. Helv. 87 (2012), no. 4, p. 1011–1033.
\bibitem[9]{duc} A. Ducros – “La structure des courbes analytiques”, https://webusers.imj-prg.fr/\textasciitilde{}antoine. ducros/trirss.pdf.
\bibitem[10]{lordan} L. Fantini \& D. Turchetti – “Triangulations of non-Archimedean curves, semi-stable reduction, and ramification”, Ann. Inst. Fourier (Grenoble) 73 (2023), no. 2, p. 695–746.
\bibitem[11]{gabber} O. Gabber, Q. Liu \& D. Lorenzini – “The index of an algebraic variety”, Invent. Math. 192 (2013), no. 3, p. 567–626.
\bibitem[12]{gillen} P. Gille \& E. Neher – “Springer’s odd degree extension theorem for quadratic forms over semilocal rings”, Indag. Math. (N.S.) 32 (2021), no. 6, p. 1290–1310.
\bibitem[13]{gruson} L. Gruson – “Une propriété des couples henséliens”, in Colloque d’algèbre commutative (Rennes, 1972), 1972, Exp. No. 10, 13 p.
\bibitem[14]{guhsur} J. Guhan – “Local-global principle for hermitian spaces over semi-global fields”, 2022, arXiv: 2203.09651.
\bibitem[15]{hhk} D. Harbater, J. Hartmann \& D. Krashen – “Applications of patching to quadratic forms and central simple algebras”, Invent. Math. 178 (2009), no. 2, p. 231–263.
\bibitem[16]{hhk10} D. Harbater, J. Hartmann \& D. Krashen – “Refinements to patching and applications to field invariants”, Internat. Math. Res. Notices (2015), no. 20, p. 10399–10450.
\bibitem[17]{harder1} G. Harder – “Über die Galoiskohomologie halbeinfacher Matrizengruppen. I”, Math. Z. 90 (1965), p. 404–428.
\bibitem[18]{harder2} G. Harder – “Über die Galoiskohomologie halbeinfacher Matrizengruppen. II”, Math. Z. 92 (1966), p. 396–415.
\bibitem[19]{harder3} G. Harder – “Über die Galoiskohomologie halbeinfacher algebraischer Gruppen. III”, J. reine angew. Math. 274/275 (1975), p. 125–138.
\bibitem[20]{dejong} A. J. de Jong – “Crystalline Dieudonné module theory via formal and rigid geometry”, Publ. Math. Inst. Hautes Études Sci. 82 (1995), p. 5–96.
\bibitem[21]{kne1} M. Kneser – “Galois-Kohomologie halbeinfacher algebraischer Gruppen über p-adischen Körpern. I”, Math. Z. 88 (1965), p. 40–47.
\bibitem[22]{kne2} M. Kneser – “Galois-Kohomologie halbeinfacher algebraischer Gruppen über p-adischen Körpern. II”, Math. Z. 89 (1965), p. 250–272.
\bibitem[23]{involutions} M.-A. Knus, A. Merkurjev, M. Rost \& J.-P. Tignol – The book of involutions, Amer. Math. Soc. Colloq. Publ., vol. 44, American Mathematical Society, Providence, RI, 1998.
\bibitem[24]{lam} T. Y. Lam – Introduction to quadratic forms over fields, Graduate Studies in Math., vol. 67, American Mathematical Society, Providence, RI, 2005.
\bibitem[25]{larmour} D. W. Larmour – “A Springer theorem for Hermitian forms”, Math. Z. 252 (2006), no. 3, p. 459– 472.
\bibitem[26]{liulibri} Q. Liu – Algebraic geometry and arithmetic curves, Oxford Graduate Texts in Math., vol. 6, Oxford University Press, Oxford, 2002.
\bibitem[27]{mart} F. Martin – “Analytic functions on tubes of nonarchimedean analytic spaces”, Algebra Number Theory 11 (2017), no. 3, p. 657–683.
\bibitem[28]{doktoratura} V. Mehmeti – “Patching on Berkovich spaces and the local–global principle”, PhD Thesis, University of Caen Normandy, 2019.
\bibitem[29]{une} V. Mehmeti – “Patching over Berkovich curves and quadratic forms”, Compositio Math. 155 (2019), no. 12, p. 2399–2438.
\bibitem[30]{muni} M. Mustaţă \& J. Nicaise – “Weight functions on non-Archimedean analytic spaces and the Kontsevich-Soibelman skeleton”, Algebraic Geom. 2 (2015), no. 3, p. 365–404.
\bibitem[31]{parimala} R. Parimala – “Homogeneous varieties—zero-cycles of degree one versus rational points”, Asian J. Math. 9 (2005), no. 2, p. 251–256.
\bibitem[32]{parsurpreeti} R. Parimala, R. Preeti \& V. Suresh – “Local-global principle for reduced norms over function fields of p-adic curves”, Compositio Math. 154 (2018), no. 2, p. 410–458.
\bibitem[33]{parsur} R. Parimala \& V. Suresh – “Local-global principle for classical groups over function fields of p-adic curves”, Comment. Math. Helv. 97 (2022), p. 255–304.
\bibitem[34]{resur} B. S. Reddy \& V. Suresh – “Admissibility of groups over function fields of p-adic curves”, Adv. Math. 237 (2013), p. 316–330.
\bibitem[35]{stacks} The Stacks Project Authors – “The Stacks Project”, 2019, https://stacks.math.columbia.edu.
\bibitem[36]{wu} Z. Wu – “Hasse principle for hermitian spaces over semi-global fields”, J. Algebra 458 (2016), p. 171–196.
\end{thebibliography}
